% =====================================================================
% CAPÍTULO 6: ALGORITMO DE SOLUCIONES RÁPIDAS Y ANALÍTICA DE DATOS
% =====================================================================
\chapter{Algoritmo de Soluciones Rápidas y Analítica de Datos}
\label{ch:algoritmo}

\section{Definición de Solución Rápida}

\begin{cajaAlgoritmo}[Definición]
Una \textbf{solución rápida} es una recomendación accionable, priorizada y
trazable que el motor analítico del \ac{SGDIC} genera para resolver una
situación detectada (saturación de cupos, queja recurrente, acumulación de
inasistencia, desempeño docente atípico o conflicto de horarios), indicando
\textbf{qué hacer, quién es el responsable, en qué plazo y con qué impacto
esperado}.
\end{cajaAlgoritmo}

\section{Arquitectura del Motor Analítico}

\begin{figure}[H]
\centering
\resizebox{0.92\textwidth}{!}{%
\begin{tikzpicture}[node distance=0.6cm and 0.8cm, every node/.style={font=\scriptsize}]

\node[terminal] (datos) {Datos operativos del \ac{SGDIC}};
\node[flujo, below=of datos] (etl) {1. Extracción y limpieza (\ac{ETL})};
\node[flujo, below=of etl] (feat) {2. Ingeniería de variables};
\node[flujo, below=of feat] (modelos) {3. Modelos analíticos};

\node[flujo, below left=1.0cm and 0.5cm of modelos] (m1) {Clasificación\\de quejas};
\node[flujo, below=1.0cm of modelos] (m2) {Predicción de\\demanda de cupos};
\node[flujo, below right=1.0cm and 0.5cm of modelos] (m3) {Preferencia de\\horarios};

\node[flujo, below=3.0cm of modelos] (prior) {4. Motor de priorización\\\scriptsize (impacto $\times$ urgencia $\div$ esfuerzo)};
\node[flujo, below=of prior] (rec) {5. Generación de\\soluciones rápidas};
\node[decision, below=of rec] (val) {¿Recomendación\\validada?};
\node[flujo, right=1.6cm of val] (ajuste) {Ajustar\\parámetros};
\node[terminal, below=of val] (tablero) {6. Tablero de\\verificación};

\draw[flecha] (datos) -- (etl);
\draw[flecha] (etl) -- (feat);
\draw[flecha] (feat) -- (modelos);
\draw[flecha] (modelos) -- (m1);
\draw[flecha] (modelos) -- (m2);
\draw[flecha] (modelos) -- (m3);
\draw[flecha] (m1) -- (prior);
\draw[flecha] (m2) -- (prior);
\draw[flecha] (m3) -- (prior);
\draw[flecha] (prior) -- (rec);
\draw[flecha] (rec) -- (val);
\draw[flecha] (val) -- node[right, font=\tiny]{Sí} (tablero);
\draw[flecha] (val) -- node[above, font=\tiny]{No} (ajuste);
\draw[flecha] (ajuste) |- (modelos);
\draw[flecha] (tablero.east) -- ++(1.6,0) |- node[near start, right, font=\tiny]{retroalimentación} (feat.east);

\end{tikzpicture}%
}
\caption{Arquitectura del motor analítico y de soluciones rápidas}
\label{fig:motor-analitico}
\end{figure}

\section{Descripción del Algoritmo}

\subsection{Vista general}
El algoritmo consta de seis etapas que transforman datos operativos en
recomendaciones priorizadas.

\begin{table}[H]
\centering
\caption{Etapas del algoritmo de soluciones rápidas}
\label{tab:etapas-algoritmo}
\small
\begin{tabularx}{\textwidth}{@{}c l X@{}}
\toprule
\rowcolor{moradoClaro}
\textbf{Etapa} & \textbf{Nombre} & \textbf{Operación} \\
\midrule
1 & Ingesta y limpieza &
    Consolidación de datos de cupos, mensajes, faltas, evaluaciones, quejas y
    cuestionarios; eliminación de duplicados y normalización. \\
2 & Ingeniería de variables &
    Cálculo de variables derivadas: tasa de ocupación, tiempo medio de
    respuesta, índice de inasistencia, sentimiento de quejas, frecuencia
    histórica de preferencia horaria. \\
3 & Modelado &
    Ejecución de modelos de clasificación, regresión y probabilidad según el
    tipo de problema detectado. \\
4 & Detección de anomalías &
    Identificación de desviaciones significativas respecto al comportamiento
    histórico esperado. \\
5 & Priorización &
    Ordenamiento de hallazgos según impacto potencial, urgencia y esfuerzo de
    implementación. \\
6 & Generación y despliegue &
    Construcción del texto de la recomendación, asignación de responsable,
    plazo sugerido y publicación en el tablero. \\
\bottomrule
\end{tabularx}
\end{table}

\subsection{Pseudocódigo del algoritmo}

\begin{lstlisting}[language=Python, caption={Algoritmo de generación de soluciones rápidas}, label={lst:algoritmo}]
ALGORITMO GenerarSolucionesRapidas(periodo):

    # Etapa 1: ingesta y limpieza
    D <- ExtraerDatos(periodo)
    D <- LimpiarYNormalizar(D)
    D <- Anonimizar(D)

    # Etapa 2: ingenieria de variables
    V <- CalcularVariables(D)

    # Etapa 3: modelado
    M_quejas  <- ModeloClasificacionQuejas(V.quejas)
    M_demanda <- ModeloPrediccionCupos(V.cupos)
    M_horario <- ModeloProbabilidadHorario(V.cuestionarios)
    M_desemp  <- ModeloDesempenoDocente(V.evaluaciones)

    # Etapa 4: deteccion de anomalias
    A <- []
    PARA CADA materia m EN V.materias HACER
        SI tasa_ocupacion(m) > 0.90 ENTONCES
            A.agregar(Anomalia("SATURACION_CUPOS", m))
        FIN SI
    FIN PARA

    PARA CADA docente d EN V.docentes HACER
        SI puntaje(d) < umbral_bajo(d) ENTONCES
            A.agregar(Anomalia("DESEMPENO_ATIPICO", d))
        FIN SI
    FIN PARA

    PARA CADA patron p EN DetectarPatrones(V.quejas) HACER
        SI p.frecuencia > umbral_frecuencia ENTONCES
            A.agregar(Anomalia("QUEJA_RECURRENTE", p))
        FIN SI
    FIN PARA

    # Etapa 5: priorizacion
    PARA CADA anomalia a EN A HACER
        a.esfuerzo <- EstimarEsfuerzo(a)
        a.score    <- (EstimarImpacto(a) * EstimarUrgencia(a))
                      / (a.esfuerzo + EPSILON)
    FIN PARA
    A <- OrdenarDescendente(A, por=a.score)

    # Etapa 6: generacion de recomendaciones
    R <- []
    PARA CADA anomalia a EN A HACER
        r <- ConstruirRecomendacion(a)
        r.responsable <- AsignarResponsable(a)
        r.plazo       <- CalcularPlazo(EstimarUrgencia(a))
        r.confianza   <- CalcularConfianza(a)
        r.estado      <- "PENDIENTE"
        R.agregar(r)
    FIN PARA

    PublicarEnTablero(R)
    RETORNAR R

FIN ALGORITMO
\end{lstlisting}

\subsection{Fórmula de priorización}

La puntuación de prioridad de cada recomendación se calcula como:

\begin{equation}
\text{Score} = \frac{I \cdot U}{E + \varepsilon}
\label{eq:score}
\end{equation}

donde $I \in [1,5]$ es el impacto estimado, $U \in [1,5]$ la urgencia,
$E \in [1,5]$ el esfuerzo de implementación y $\varepsilon$ un valor pequeño
que evita la división por cero. Este criterio favorece las acciones de alto
impacto y alta urgencia con bajo esfuerzo, es decir, las soluciones rápidas
propiamente dichas.

\subsection{Modelo de preferencia de horarios por probabilidad}

El módulo de cuestionarios de materias y horarios alimenta el modelo
probabilístico de preferencia horaria. Para cada combinación de perfil de
estudiante $p$ y franja horaria $h_i$, se estima la probabilidad:

\begin{equation}
P(h_i \mid p) =
\frac{P(p \mid h_i) \cdot P(h_i)}{P(p)}
= \frac{N_{p,h_i} + \alpha}{N_{h_i} + \alpha \cdot |P|}
\label{eq:prob-horario}
\end{equation}

donde $N_{p,h_i}$ es el número de estudiantes del perfil $p$ que prefirieron la
franja $h_i$, $N_{h_i}$ el total de preferencias de esa franja, $|P|$ el número
de perfiles y $\alpha$ un parámetro de suavizado de Laplace que evita
probabilidades nulas para franjas sin observaciones.

La preferencia esperada total de una franja se calcula como el promedio
ponderado sobre los perfiles de la población objetivo:

\begin{equation}
P(h_i) = \sum_{p \in P} P(h_i \mid p) \cdot w_p
\label{eq:prob-total}
\end{equation}

donde $w_p$ es el peso relativo del perfil $p$ en la matrícula. La asignación de
horario se resuelve como un problema de optimización que maximiza la aceptación
esperada:

\begin{equation}
\max \sum_{i=1}^{n} P(h_i) \cdot x_i
\quad \text{sujeto a} \quad \sum_{i=1}^{n} x_i = 1, \; x_i \in \{0,1\}
\label{eq:optimizacion}
\end{equation}

donde $x_i = 1$ indica que la franja $h_i$ es seleccionada para la materia.

\section{Modelos Analíticos por Tipo de Problema}

\begin{table}[H]
\centering
\caption{Modelos analíticos aplicados por tipo de problema}
\label{tab:modelos}
\small
\begin{tabularx}{\textwidth}{@{}p{3.2cm} p{6.6cm} X@{}}
\toprule
\rowcolor{azulClaro}
\textbf{Problema} & \textbf{Modelo} & \textbf{Salida} \\
\midrule
Clasificación de quejas &
  Clasificación supervisada (Naive Bayes / \ac{KNN}) sobre el texto de la queja &
  Categoría, área responsable y urgencia sugerida \\
Predicción de demanda de cupos &
  Series temporales (suavizado exponencial / ARIMA) por materia y periodo &
  Demanda esperada y riesgo de saturación \\
Preferencia de horarios &
  Probabilidad condicional con suavizado de Laplace &
  $P(h_i \mid p)$ y ranking de franjas preferidas \\
Desempeño docente atípico &
  Control estadístico de procesos ($\mu \pm 2\sigma$) &
  Docentes con puntaje fuera del rango esperado \\
Reincidencia de inasistencia &
  Regresión logística sobre historial de faltas &
  Probabilidad de superar el umbral permitido \\
Tiempo de respuesta &
  Estadística descriptiva y percentiles &
  Materias y actores con tiempos sobre el objetivo \\
Temas recurrentes en quejas &
  Procesamiento de lenguaje natural (frecuencia de términos) &
  Temas emergentes y su evolución temporal \\
\bottomrule
\end{tabularx}
\end{table}

\section{Reglas de Generación de Soluciones Rápidas}

\begin{table}[H]
\centering
\caption{Catálogo de soluciones rápidas por anomalía detectada}
\label{tab:catalogo-soluciones}
\small
\begin{tabularx}{\textwidth}{@{}p{3.2cm} X p{3.0cm}@{}}
\toprule
\rowcolor{verdeClaro}
\textbf{Anomalía} & \textbf{Solución rápida sugerida} & \textbf{Responsable} \\
\midrule
Saturación de cupos ($>90\%$) &
  Abrir una sección adicional o ampliar el cupo máximo; notificar a los
  estudiantes en lista de espera con horario alternativo. & Departamento \\
Demanda proyectada creciente &
  Planificar apertura de sección antes del inicio del periodo de matrícula. &
  Departamento \\
Queja recurrente en una materia &
  Reunión formal con el docente, revisión del plan de asignatura y plan de
  mejoramiento con seguimiento quincenal. & Secretaría + Departamento \\
Queja crítica sin atender ($>3$ días) &
  Escalamiento inmediato a dirección y notificación al responsable con plazo
  de 24 horas. & Dirección \\
Desempeño docente atípico &
  Activar plan de acompañamiento pedagógico y observación de clase. &
  Departamento \\
Inasistencia acumulada $>20\%$ &
  Alerta al estudiante con plan de recuperación y citación a tutoría. &
  Docente + Secretaría \\
Evaluación docente incompleta &
  Recordatorio intensivo y habilitación de tiempo protegido en clase. &
  Docente + Sistema \\
Conflicto de horarios frecuente &
  Reubicar la materia a la franja con mayor $P(h_i)$ estimada. & Departamento \\
Mensajes sin respuesta $>48$ h &
  Recordatorio automático al destinatario y copia a la secretaría. & Sistema \\
Tiempo de respuesta a cupos elevado &
  Redistribuir la carga de revisión entre el personal de secretaría. &
  Secretaría \\
\bottomrule
\end{tabularx}
\end{table}

\section{Tablero de Verificación de Resultados}

\begin{figure}[H]
\centering
\begin{tikzpicture}[node distance=0.45cm and 0.5cm, every node/.style={font=\scriptsize}]

\node[flujo, minimum width=3.4cm] (tab) at (0,0) {Tablero de\\Verificación};

\node[flujo, below left=1.0cm and 0.3cm of tab] (t1) {Resultados\\académicos};
\node[flujo, below=1.0cm of tab] (t2) {Quejas y\\sugerencias};
\node[flujo, below right=1.0cm and 0.3cm of tab] (t3) {Evaluación\\docente};

\node[flujo, below=1.0cm of t1] (v1) {Ocupación de cupos\\Tiempo de respuesta};
\node[flujo, below=1.0cm of t2] (v2) {Tasa de cierre\\Tiempo por caso};
\node[flujo, below=1.0cm of t3] (v3) {Participación\\Puntaje promedio};

\node[decision, below=1.0cm of v2] (dec) {¿Meta\\alcanzada?};
\node[terminal, left=1.6cm of dec] (ok) {Reportar\\logro};
\node[flujo, right=1.6cm of dec] (no) {Generar acción\\correctiva};

\draw[flecha] (tab) -- (t1);
\draw[flecha] (tab) -- (t2);
\draw[flecha] (tab) -- (t3);
\draw[flecha] (t1) -- (v1);
\draw[flecha] (t2) -- (v2);
\draw[flecha] (t3) -- (v3);
\draw[flecha] (v1) |- (dec);
\draw[flecha] (v3) |- (dec);
\draw[flecha] (v2) -- (dec);
\draw[flecha] (dec) -- node[above, font=\tiny]{Sí} (ok);
\draw[flecha] (dec) -- node[above, font=\tiny]{No} (no);
\draw[flecha] (no) |- (tab);

\end{tikzpicture}
\caption{Estructura del tablero de verificación de resultados}
\label{fig:tablero}
\end{figure}

\subsection{Elementos del tablero}
\begin{table}[H]
\centering
\caption{Componentes del tablero de verificación}
\label{tab:tablero}
\small
\begin{tabularx}{\textwidth}{@{}p{3.3cm} X p{3.1cm}@{}}
\toprule
\rowcolor{moradoClaro}
\textbf{Sección} & \textbf{Contenido} & \textbf{Usuario} \\
\midrule
Panel de cupos & Ocupación por materia, lista de espera, tiempo de respuesta. & Secretaría, Departamento \\
Panel de quejas & Casos abiertos, vencidos, por categoría y responsable. & Secretaría, Dirección \\
Panel de evaluación & Participación por materia, puntajes agregados, alertas. & Departamento, Docente \\
Panel de comunicación & Mensajes sin respuesta, tiempos promedio, escalamientos. & Secretaría \\
Panel de faltas & Estudiantes en riesgo, alertas por umbral, justificaciones pendientes. & Docente, Secretaría \\
Panel de soluciones & Recomendaciones activas, aplicadas y su efecto medido. & Departamento \\
\bottomrule
\end{tabularx}
\end{table}

\section{Métricas de Evaluación del Algoritmo}

\begin{table}[H]
\centering
\caption{Métricas de desempeño del motor analítico}
\label{tab:metricas-algoritmo}
\small
\begin{tabularx}{\textwidth}{@{}p{3.6cm} X c@{}}
\toprule
\rowcolor{azulClaro}
\textbf{Métrica} & \textbf{Definición} & \textbf{Meta} \\
\midrule
Exactitud & $\frac{TP+TN}{TP+TN+FP+FN}$ & $> 0.85$ \\
Precisión & $\frac{TP}{TP+FP}$ & $> 0.80$ \\
Exhaustividad & $\frac{TP}{TP+FN}$ & $> 0.80$ \\
$F_1$ & Media armónica de precisión y exhaustividad & $> 0.80$ \\
Error absoluto medio & $\frac{1}{n}\sum |y_i - \hat{y}_i|$ (demanda de cupos) & $< 10\%$ \\
Tasa de aceptación & Recomendaciones aplicadas / generadas & $> 70\%$ \\
Efectividad & Mejora medida tras la aplicación & $> 60\%$ \\
Cobertura & Anomalías detectadas / anomalías reales & $> 85\%$ \\
\bottomrule
\end{tabularx}
\end{table}

\section{Ejemplo de Aplicación}

\begin{cajaAlgoritmo}[Ejemplo: quejas recurrentes en una materia de programación]
\textbf{Datos observados:} 18 quejas en el periodo correspondiente a la materia
\textit{Programación II}, 14 de ellas con el término ``avance'' y ``ritmo'';
la tasa de inasistencia del grupo es del 24\%; el tiempo medio de respuesta del
docente es de 96 horas.

\vspace{4pt}
\textbf{Hallazgos del motor:}
\begin{itemize}[leftmargin=1.0cm, itemsep=2pt]
  \item Anomalía \texttt{QUEJA\_RECURRENTE} (frecuencia $>$ umbral).
  \item Anomalía \texttt{INASISTENCIA\_ALTA} ($24\% > 20\%$).
  \item Anomalía \texttt{RESPUESTA\_LENTA} ($96$ h $> 48$ h).
\end{itemize}

\vspace{4pt}
\textbf{Priorización:} impacto $I=5$, urgencia $U=4$, esfuerzo $E=2$
$\Rightarrow$ Score $= (5 \times 4)/(2 + 0.01) \approx 9.95$ (prioridad alta).

\vspace{4pt}
\textbf{Solución rápida generada:} reunión formal entre secretaría, docente y
coordinación académica en un plazo de 3 días; revisión del cronograma del plan
de asignatura; tutoría de refuerzo semanal; recordatorios automáticos al docente
sobre mensajes pendientes.
\end{cajaAlgoritmo}